\documentclass[instructions.tex]{subfiles}
\begin{document}
 
\chapter{De la rechute dans le péché mortel}
 
\primo Quand on retombe dans le péché mortel, rarement ou dans des occasions imprévues, et qu’aussitôt on s’en repent, qu’on recourt à Dieu par le sacrement de pénitence et par la prière, on peut présumer que ces rechutes sont arrivées par surprise et par fragilité. De telles chutes doivent faire sentir à une âme qu’elle n’est que foiblesse, qu’elle doit se défier de ses forces, et travailler à son salut avec tremblement.
 
Mais si, faute d’éviter les occasions et de veiller sur soi, on retombe souvent dans les péchés mortels, c’est alors par une négligence crasse ou par malice. O Dieu ! que l’état d’une âme qui tombe ainsi habituellement dans le péché mortel, est déplorable ! Quand on n’aurait commis qu’un seul péché mortel en sa vie, c’est assez, dit Tertullien, pour pleurer le reste de ses jours, et pour pleurer toute l’éternité ; mais, loin de pleurer les péchés passés, vous en ajoutez de nouveaux.
 
Combien de gens plus foibles que vous, qui ont des tentations plus fréquentes, et moins de grâces, et qui ne retombent point ! Pourquoi n’avez-vous pas le courage de les imiter ? Il vaudrait mieux perdre la vie que de perdre la grâce de Dieu ; néanmoins vous la perdez avec autant de facilité que si vous ne perdiez qu’une obole : perte d’autant plus funeste, que vous ne la comprenez pas.
 
Les tyrans employaient les supplices pour faire tomber un chrétien dans une infidélité ; et vous, qu’est-ce qui vous fait tomber ? Ce ne sont ni des bourreaux, ni des tourmens ; c’est une tentation, une amitié, une fréquentation qu’il ne vous plaît pas de quitter, une attache aux plaisirs de votre corps ou à vos intérêts. O lâcheté criminelle !
 
\secundo Craignez qu’en multipliant vos péchés et vos rechutes, Dieu ne vous y laisse ! Que savez-vous si le premier péché mortel que vous commettrez, ne mettra point le sceau à votre réprobation ? Il est vrai que Dieu pardonne toutes les fois qu’on se convertit sincérement ; mais en aurez-vous à l’avenir la volonté ? Il est à craindre que vous ne l’ayez pas, parce que plus vous retombez, plus votre cœur s’endurcit, plus il est difficile de vous convertir de nouveau. D’ailleurs, plus vous multipliez vos rechutes, plus Dieu s’éloigne de vous ; sa grâce devient plus rare, vous devenez plus foible, et le démon plus puissant. En vous convertissant, vous aviez chassé l’ennemi de votre cœur ; si vous retombez, il y rentrera, dit Jésus-Christ, avec sept autres démons plus méchans.
 
Profitez de cet avis de saint Bernard : Craignez, mon frère, quand vous avez reçu la grâce et l’amitié de Dieu ; craignez bien plus, quand vous l’avez perdue ; mais craignez beaucoup plus, quand vous l’avez recouvrée, parce que, si vous la perdez une seconde et une troisième fois, peut-être ne la recouvrerez-vous jamais\footnote{Time, frater, pro accepta gratia, ampliùs pro amissâ, longè plus pro recuperatâ. S. Bern.}
 
\section{Résolutions}
 
\primo Je me pénétrerai de plus en plus d’horreur pour la rechute dans le péché mortel.
 
\secundo A cet effet, je relirai souvent ce chapitre.
 
\tertio Et je prendrai avec mon guide spirituel les mesures convenables pour m’aider à éviter désormais ce malheur. 

\prayer{O mon Dieu ! c’est de vous que j’attends les secours dont j’ai besoin pour me convertir enfin solidement. Ayez pitié de moi ; sauvez-moi.}
 
\end{document}
